\scanpage{6}
\mspage{4}
Donc on aurait envie de prendre
\[
\deleted{\mathbb Q}\ \mathbb Z\big(\{\unsure{(}X_i\unsure{)}\}[dX_i]\big)_{0\leq i\leq n}/\big(\textstyle\sum X_i-1,\sum dX_i\big)_{\deleted{\mathrm{pd}}}
\]
où $\{\ \}$ désigne les pol. à p.d., mais ce n'est plus $\simeq\mathbb Z\{X_0,\ldots,X_{n-1}\}[dX_0,\ldots,dX_{n-1}]$ \deleted{(qui est} \deleted{\unsure{le est|la cat}}\inserted{bien} une résolution de $\mathbb Z$), car si on veut tirer de la relation $\sum_0^n X_i-1=0$,
\[
X_n=1-\sum_0^{n-1}X_i\quad(\inserted{\text{et}}\ dX_n=-\sum_0^{n-1}dX_i\quad\text{de }\sum_{\unsure{1|0}}^n dX_i=0)\ .\text{--}
\]
on a le \unsure{``bec''|``hic''} que $1-\sum_0^{n-1}X_i$ n'appartient pas à l'idéal à p.d. donné dans $\mathbb Z\{X_0,\ldots,X_{n-1}\}$ (donc on ne voit pas comment envoyer $\mathbb Z\{X_0,\ldots,X_{n-1}\}$ dans $\mathbb Z\{X_0,\ldots,X_{n-1}\}$ avec \unsure{comme|même} \deleted{noyau} $\sum_{\unsure{1|0}}^n X_i-1$ dans le noyau\ldots).

On s'en tire en prenant un anneau de coeff. différent de $\mathbb Z$, soit $S$, avec un \deleted{indéterminé}\inserted{``paramètre''} $t\in S$ fixé dont on \uline{sache prendre les p.d.} \inserted{(i.e. $t\in J$, $J$ idéal à p.d. de $S$)}, et en remplaçant les équations $\sum X_i=1$ de $E^{[n]}$ par l'équation
\[
\sum X_i=t\qquad\text{dans } S^{n+1}
\]
et définissant
\[
\big(C^\bullet_{\mathrm{DRpd}\,[n]}(S,J,t)\big)\unsure{\simeq|=}
\big(S\{X_i\}_{0\leq i\leq n}[dX_i]_{0\leq i\leq n}\big)/
\big(\textstyle\sum X_i-t,\sum dX_i\big)
\]
(On divise par l'idéal \uline{à p.d.} engendré \unsure{pd} par $\sum X_i-t,\ \sum dX_i$, ce qui a un sens
% archive sheet number: 5
